% GENERATED FILE. Edit canonical lessons, then run npm run generate:books.
% canonical-source-hash: fnv1a64:35646e21f8f72bc2
% canonical-lessons: PT-C458-serra, PT-C458-encosta, PT-C458-cume, PT-C458-desfiladeiro, PT-C458-escalar

\chapter{Na serra}
\label{ch:pt-na-serra}
\hlchaptermodality{458}

\begin{chapteropening}
\textbf{By the end of this chapter:} \emph{I can describe a mountain walk: the mountain range, the hillside, the summit, a gorge and climbing.}

Complete the last lesson of chapter 458: I can describe a mountain walk: the mountain range, the hillside, the summit, a gorge and climbing.
\end{chapteropening}

\section[a serra]{a serra — the mountain range, the hills; the saw}

\label{lesson:PT-C458-serra}



\begin{quote}
Before the new one: say the Portuguese for cork, then the Portuguese for to prune.
\end{quote}

\subsection*{You'll want to know: a serra}
\textbf{a serra} — \textquotedblleft{}the mountain range, the hills; the saw\textquotedblright{}.

\begin{quote}
Passámos o fim de semana na serra.
\end{quote}

\begin{quote}
\emph{We spent the weekend in the mountains.}
\end{quote}

A ridge of peaks looks like the teeth of a saw. \emph{Serrar} — to saw.

\begin{cousinweb}
\emph{serra} is inherited from Latin \emph{serra} 'saw' (EtymDB): a mountain ridge has the jagged outline of a saw's teeth.
\end{cousinweb}

\begin{grammarlens}[title={What you've built}]
One more word for telling, explaining and describing at B1.
\end{grammarlens}

\subsection*{Your turn}
\begin{itemize}
\raggedright
  \item \emph{Say it:} \emph{a serra}
  \item \emph{Say it:} \emph{a serra}, once more
  \item \emph{Say it:} a sentence of your own with \emph{a serra}
  \item \emph{Recall:} say the Portuguese for cork, then the Portuguese for to prune, then say \emph{a serra} again
\end{itemize}

\begin{tcolorbox}[breakable,colback=teal!4,colframe=teal!35!black,title={Before you move on}]
What does a serra mean? (\textquotedblleft{}The mountain range, the hills; the saw\textquotedblright{}.) Say it once more.
\end{tcolorbox}
\section[a encosta]{a encosta — the hillside, the slope}

\label{lesson:PT-C458-encosta}



\begin{quote}
Before the new one: say the Portuguese for to prune, then the Portuguese for the mountain range.
\end{quote}

\subsection*{You'll want to know: a encosta}
\textbf{a encosta} — \textquotedblleft{}the hillside, the slope\textquotedblright{}.

\begin{quote}
Há vinhas na encosta do rio.
\end{quote}

\begin{quote}
\emph{There are vineyards on the slope above the river.}
\end{quote}

\emph{Encosta abaixo} — downhill.

\begin{cousinweb}
\emph{encosta} is built on \emph{encostar} 'to lean against' (EtymDB), on \emph{costa} 'side, rib', from Latin \emph{costa} 'rib' (Lewis \& Short).
\end{cousinweb}

\begin{grammarlens}[title={What you've built}]
One more word for telling, explaining and describing at B1.
\end{grammarlens}

\subsection*{Your turn}
\begin{itemize}
\raggedright
  \item \emph{Say it:} \emph{a encosta}
  \item \emph{Say it:} \emph{a encosta}, once more
  \item \emph{Say it:} a sentence of your own with \emph{a encosta}
  \item \emph{Recall:} say the Portuguese for to prune, then the Portuguese for the mountain range, then say \emph{a encosta} again
\end{itemize}

\begin{tcolorbox}[breakable,colback=teal!4,colframe=teal!35!black,title={Before you move on}]
What does a encosta mean? (\textquotedblleft{}The hillside, the slope\textquotedblright{}.) Say it once more.
\end{tcolorbox}
\section[o cume]{o cume — the summit, the top}

\label{lesson:PT-C458-cume}



\begin{quote}
Before the new one: say the Portuguese for the mountain range, then the Portuguese for the hillside.
\end{quote}

\subsection*{You'll want to know: o cume}
\textbf{o cume} — \textquotedblleft{}the summit, the top\textquotedblright{}.

\begin{quote}
Chegámos ao cume ao meio-dia.
\end{quote}

\begin{quote}
\emph{We reached the summit at midday.}
\end{quote}

Also the peak of a career.

\begin{cousinweb}
\emph{cume} goes back to Latin \emph{culmen} 'top, roof' (Lewis \& Short); EtymDB lists the Portuguese word without a source.
\end{cousinweb}

\begin{grammarlens}[title={What you've built}]
One more word for telling, explaining and describing at B1.
\end{grammarlens}

\subsection*{Your turn}
\begin{itemize}
\raggedright
  \item \emph{Say it:} \emph{o cume}
  \item \emph{Say it:} \emph{o cume}, once more
  \item \emph{Say it:} a sentence of your own with \emph{o cume}
  \item \emph{Recall:} say the Portuguese for the mountain range, then the Portuguese for the hillside, then say \emph{o cume} again
\end{itemize}

\begin{tcolorbox}[breakable,colback=teal!4,colframe=teal!35!black,title={Before you move on}]
What does o cume mean? (\textquotedblleft{}The summit, the top\textquotedblright{}.) Say it once more.
\end{tcolorbox}
\section[o desfiladeiro]{o desfiladeiro — the gorge, the narrow pass}

\label{lesson:PT-C458-desfiladeiro}



\begin{quote}
Before the new one: say the Portuguese for the hillside, then the Portuguese for the summit.
\end{quote}

\subsection*{You'll want to know: o desfiladeiro}
\textbf{o desfiladeiro} — \textquotedblleft{}the gorge, the narrow pass\textquotedblright{}.

\begin{quote}
O rio corre no fundo do desfiladeiro.
\end{quote}

\begin{quote}
\emph{The river runs at the bottom of the gorge.}
\end{quote}

A pass so narrow that people walk in single file.

\begin{cousinweb}
\emph{desfiladeiro} is built on \emph{desfilar} 'to march in file', which was borrowed from French \emph{défiler}, on \emph{fil} 'thread', Latin \emph{fīlum} (Lewis \& Short).
\end{cousinweb}

\begin{grammarlens}[title={What you've built}]
One more word for telling, explaining and describing at B1.
\end{grammarlens}

\subsection*{Your turn}
\begin{itemize}
\raggedright
  \item \emph{Say it:} \emph{o desfiladeiro}
  \item \emph{Say it:} \emph{o desfiladeiro}, once more
  \item \emph{Say it:} a sentence of your own with \emph{o desfiladeiro}
  \item \emph{Recall:} say the Portuguese for the hillside, then the Portuguese for the summit, then say \emph{o desfiladeiro} again
\end{itemize}

\begin{tcolorbox}[breakable,colback=teal!4,colframe=teal!35!black,title={Before you move on}]
What does o desfiladeiro mean? (\textquotedblleft{}The gorge, the narrow pass\textquotedblright{}.) Say it once more.
\end{tcolorbox}
\section[escalar]{escalar — to climb (a rock or a mountain)}

\label{lesson:PT-C458-escalar}



\begin{quote}
Before the new one: say the Portuguese for the summit, then the Portuguese for the gorge.
\end{quote}

\subsection*{You'll want to know: escalar}
\textbf{escalar} — \textquotedblleft{}to climb (a rock or a mountain)\textquotedblright{}.

\begin{quote}
Eles escalaram a rocha sem medo.
\end{quote}

\begin{quote}
\emph{They climbed the rock without fear.}
\end{quote}

\emph{Escalar uma montanha} — to climb a mountain. In Brazil, \emph{escalar} also means to pick players for a team.

\begin{cousinweb}
\emph{escalar} is built from \emph{escala} 'ladder' and \emph{-ar} (EtymDB); \emph{escala} is a learned borrowing from Latin \emph{scāla} 'flight of steps'.
\end{cousinweb}

\begin{grammarlens}[title={What you've built}]
One more word for telling, explaining and describing at B1.
\end{grammarlens}

\subsection*{Your turn}
\begin{itemize}
\raggedright
  \item \emph{Say it:} \emph{escalar}
  \item \emph{Say it:} \emph{escalar}, once more
  \item \emph{Say it:} a sentence of your own with \emph{escalar}
  \item \emph{Recall:} say the Portuguese for the summit, then the Portuguese for the gorge, then say \emph{escalar} again
\end{itemize}

\begin{tcolorbox}[breakable,colback=teal!4,colframe=teal!35!black,title={Before you move on}]
What does escalar mean? (\textquotedblleft{}To climb (a rock or a mountain)\textquotedblright{}.) Say it once more.
\end{tcolorbox}
